\documentclass{article}
\usepackage[utf8]{inputenc}
\usepackage[T1]{fontenc}
\usepackage[a4paper,total={6in,8in}]{geometry}
\usepackage{amsmath,amssymb,amsfonts}
\usepackage{graphicx}
\usepackage{hyperref}
\usepackage{enumitem}
\usepackage{listings}
\usepackage{xcolor}
\DeclareMathOperator{\Var}{Var}
\DeclareMathOperator{\Cov}{Cov}
\providecommand{\tightlist}{\setlength{\itemsep}{0pt}\setlength{\parskip}{0pt}}
\lstset{basicstyle=\ttfamily\scriptsize,breaklines=true,columns=fullflexible,
  keepspaces=true,showstringspaces=false,frame=single}
\hypersetup{colorlinks=true,urlcolor=blue,linkcolor=blue}
\setlength{\emergencystretch}{2em}

\title{Math 126 Problem 7: Deriving the Normal Distribution}
\author{} % Put your name here.
\date{Spring 2027}
\begin{document}
\maketitle

\hypertarget{submission-instructions}{%
\subsection*{Submission Instructions}\label{submission-instructions}}

Please upload your work to Gradescope by \textbf{11:59pm on Monday, February 15, 2027 (Pacific time).}

\begin{itemize}
\item
  You are encouraged to discuss ideas and work with your classmates. However, you \emph{must acknowledge} your collaborators at the top of each solution on which you collaborated with others and you \emph{must write} your solutions independently.
\item
  Your solutions to theory questions must be written legibly, or typeset in LaTeX or markdown.
\item
  I recommend that you write your solutions in a Python notebook using Google Colab.
\item
  You should submit your solutions as a \emph{single PDF} via the assignment on Gradescope.
\end{itemize}

\textbf{Grading:} The point of the problem set is for \emph{you} to learn. To this end, I hope to disincentivize the use of LLMs by \emph{not} grading your work for correctness. Instead, you will grade your own work by comparing it to my solutions. This self-grade is due the Friday \emph{after} the problem set is due, also on Gradescope.

\noindent Original notebook: \href{https://colab.research.google.com/drive/1ciQVopHz3Cb9xFEL5UnwU49hLGCGxnF_?usp=sharing}{Google Colab}.

% Source notebook cell 2
\hypertarget{deriving-the-normal-distribution}{%
\subsection*{Deriving the Normal Distribution}\label{deriving-the-normal-distribution}}

% Source notebook cell 3
\begin{center}
\includegraphics[width=0.65\linewidth,height=0.28\textheight,keepaspectratio]{images/problem07-figure1.jpg}
\end{center}

% Source notebook cell 4
In this problem, we will derive the density function of the normal distribution from the example of a dartboard. Let \(f: \mathbb{R}^d \to [0,1]\) be a probability density function that describes the probability a dart lands at the point \((x,y)\). We want our probability density function \(f\) to have two properties:

\begin{itemize}
\item
  \textbf{Radial Symmetry:} The probability that a dart lands at a point depends only on the distance between the point and the origin.
\item
  \textbf{Independence:} Coordinates are independent e.g., knowing the \(x\)-coordinate does not give us information about the \(y\)-coordinate.
\end{itemize}

From the radial symmetry property, we can conclude that \begin{align*}
f(x,y) = f(r) = f(\sqrt{x^2 + y^2})
\end{align*} where \(r=\sqrt{x^2+y^2}\) is the distance between the origin and \((x,y)\).

From the independence property, we can conclude that \(f(x,y) = g(x) h(y)\) for some functions \(g\) and \(h\). Further, the radial symmetry property tells us that \(f(x,y) = f(y,x) = g(y) h(x)\) so \(g\) and \(h\) must be the same function. That is, \begin{align*}
f(x,y) = g(x) g(y).
\end{align*}

% Source notebook cell 5
\hypertarget{part-a}{%
\subsection*{Part A}\label{part-a}}

Up to rescaling, we can assume that \(g(0)=1\). Use this fact and plug in the point \((r,0)\) to conclude that \(f\) and \(g\) are the same function.

% Source notebook cell 6
\hypertarget{part-b}{%
\subsection*{Part B}\label{part-b}}

Define the function \(h(x) = f(\sqrt{x})\). In particular, \begin{align*}
f(\sqrt{x^2 + y^2}) = f(x) f(y) \Leftrightarrow
h(x^2 + y^2) = h(x^2) h(y^2).
\end{align*} Use this property to show that \begin{align*}
h(x_1 + x_2 + \ldots + x_n) = h(x_1) h(x_2) \cdot \ldots \cdot h(x_n).
\end{align*} Then prove that \(h(n) = h(1)^n\).

Without loss of generality, let \(h(1) = b\).

% Source notebook cell 7
\hypertarget{part-c}{%
\subsection*{Part C}\label{part-c}}

Let \(p\) and \(q\) be any integers. Prove that \begin{align*}
h\left(\frac{p}{q}\right) = b^{\frac{p}{q}}.
\end{align*}

\textbf{Hint:} Consider the expression \begin{align*}
h\left(\frac{p}{q} \cdot q \right).
\end{align*}

As long as \(h\) is continuous, conclude that \(h(x) = b^x\) for any real number since the rational numbers are dense in the real number line.

% Source notebook cell 8
\hypertarget{part-d}{%
\subsection*{Part D}\label{part-d}}

You have shown that \(h(x) = b^x\). Further, \(h(x) = e^{cx}\) for some real number \(c\). Then \(f(x) = h(x^2) = e^{c x^2}\). Recall this implies the density function of \(x\) and \(y\) is \(f(x,y) = f(x) f(y) = e^{c (x^2 + y^2)}\).

For \(f\) to be a valid probability density function, what constraint do we have on \(c\)? If we wanted to make the distribution more concentrated, how should we change \(c\)?

Except for the normalization, we have explained every part of the multivariate normal distribution given below \begin{align*}
f(x,y) = \frac1{\sigma^2 \cdot 2 \pi} \exp\left(-\frac{x^2 + y^2}{2\sigma^2}\right).
\end{align*} Why is there a normalization?

\end{document}
